\subsection{Summary of the work completed}
A handwriting robot was designed and built to recognise photographed text, identify a known writer and reproduce new digital text in that writer's style. The system combines page conditioning, line segmentation, neural recognition, stored glyph trajectories and a calibrated Cartesian gantry. Eleven writer profiles are stored and exercised in physical reproduction.

\subsection{Summary of observations and findings}
The selected reader achieves 94.06\% character accuracy and writer identification reaches 98.08\% on the recorded validation set. Physical reproduction gives 84.41\% read-back character accuracy, with seven writers above 85\%. Writer identity is retained on 86.41\% of reproduced lines and all fourteen page pools. The measured shape correlation is 0.697, showing that recognisable identity does not imply complete geometric similarity.

All 25 completion-confirmed ODROID reproduction jobs satisfy 4 s per character. The two measured classification requests exceed their time allowances. Every measured square and triangle length is within 0.5 mm of nominal, with a maximum difference of 0.37 mm. These results establish dimensional performance on the specimens, while absolute placement and repeatability need further qualification.

\subsection{Contribution}
The project demonstrates a complete path from offline handwriting samples to physical multi-writer output. Its first-principles network implementation is supported by numerical checks and trained models with scores close to the reference route. Explicit glyphs and G-code make the relationship between image features and pen movement visible, allowing each stage to be inspected and improved.

\subsection{Suggestions for future work}
The first priority is faster embedded classification without changing the detected lines or decoded text. Batching and compiled numerical kernels should be measured on the ODROID. Context-specific glyphs and multi-character units should then be tested on the more connected writers. Repeated bed-coordinate measurements, word overlays and matched original-to-reproduced sentences would strengthen the mechanical and stylistic evaluation. Improvements to motor protection and pulse scheduling would complete the next hardware revision.
